\section{Matched RTX Target-Only Bridge}
\label{app:bridge}
\paragraph{Design and exposure.}
The 19-task bridge is a separate fresh experiment following a retrospective comparison of already exposed outputs. All task IDs, prompt bytes/token IDs, conditions, run order, complete-output rules and retry policy were frozen before formal execution. It uses Qwen3-8B/0.6B Q8\_0, the same b10839 CUDA numerical libraries, context 4096, batch capacity 2048, F16 KV and flash attention off. Target offload is 37/37 layers and native draft offload 29/29. Numerical thread settings are 24/24; HTTP housekeeping threads are separate. N0/N1 retain target prefill UB2, draft UB2 and $k=2$, with an unchanged-capacity sham for N0 and target-continuation UB1 for N1. U2 is the ordinary released target-only CLI with global UB2.

The first repetition orders N0/N1/U2, the second N1/U2/N0 and the third U2/N0/N1. A single GPU queue uses fresh processes, no warmup and no KV reuse; OS file cache is not flushed. This is not a performance protocol. Twelve engineering generations use /15, /25, /18 and /31. The unchanged capture frontend reuses its prior 120-array/30-sham qualification, with local input/configuration/handoff checks. All 171 formal generations succeed on their first valid attempt. No additional diagnostic generation is needed.

\paragraph{Endpoints and every task.}
Arrays include an earlier EOG or all 128 committed IDs at the cap. Raw logs contain no committed overshoot in this bridge. Sixty native records retain one additional uncommitted target sample each, outside output equality. The original endpoint field is the termination reason; the added strict boundary includes length and terminal ID. The complete-array primary result is unchanged by this clarification. Table~\ref{tab:bridgeall} lists every task; lengths are identical across repetitions, and each displayed comparison holds in all three repetitions.
\begin{table*}[t]
\centering
\caption{All 19 bridge tasks. Lengths include the committed terminal ID. R denotes the preserved historical reference. Full arrays and boundaries match across all three repetitions within each condition.}
\label{tab:bridgeall}
\begin{tabular}{@{}lrrrllll@{}}
\toprule
HumanEval & $|N_0|$ & $|N_1|$ & $|U_2|$ & Stop & $N_0=R$ & $N_1=R$ & $U_2=N_1$ \\
\midrule
/143 & 88 & 88 & 88 & EOS & yes & yes & yes \\
/31 & 128 & 128 & 128 & limit & yes & yes & yes \\
/25 & 75 & 75 & 75 & EOS & yes & no & yes \\
/81 & 128 & 128 & 128 & limit & yes & yes & yes \\
/15 & 25 & 28 & 28 & EOS & no & yes & yes \\
/141 & 102 & 102 & 102 & EOS & yes & yes & yes \\
/142 & 64 & 64 & 64 & EOS & yes & yes & yes \\
/18 & 69 & 69 & 69 & EOS & no & no & yes \\
/136 & 72 & 72 & 72 & EOS & yes & yes & yes \\
/99 & 58 & 58 & 58 & EOS & yes & yes & yes \\
/148 & 112 & 112 & 112 & EOS & yes & yes & yes \\
/140 & 112 & 112 & 112 & EOS & yes & yes & yes \\
/87 & 64 & 64 & 64 & EOS & yes & yes & yes \\
/76 & 59 & 59 & 59 & EOS & yes & yes & yes \\
/26 & 114 & 114 & 114 & EOS & yes & yes & yes \\
/49 & 128 & 128 & 128 & limit & yes & yes & yes \\
/40 & 108 & 108 & 108 & EOS & yes & yes & yes \\
/109 & 66 & 66 & 66 & EOS & yes & yes & yes \\
/66 & 38 & 38 & 38 & EOS & yes & yes & yes \\
\bottomrule
\end{tabular}
\end{table*}


\paragraph{Aligned witnesses.}
On /15, /25, /18 and /31, N0/N1 initial source KV and pending logits are byte-equal. This checks the first-pending handoff, not the later first-divergence state. At /15's first differing output, index 15, N0 selects ID 9147 (score 50.415562) over 4199 (50.238899); N1 selects 4199 (50.456791) over 9147 (50.248737). Both native caller lengths are three with positions [103,104,105] and logits index two. At /25, index four, N0 selects 1406 (44.883926) over 271 (44.487068); N1 selects 271 (44.887604) over 1406 (44.698635), with positions [154,155,156] and logits index zero. These are recorded numerical witnesses, not a calibrated predictor or a causal proof.

Native target continuation observes occupied tokens [1,2] for N0 and [1] for N1; draft continuation retains [1,2]. U2 post-first-commit batch dumps have length one, while its first prefill dump spans the entire prompt and additional split dumps have one/two tokens. Caller dumps are not native tensor-callback observations. Live U2 graph tensors/KV/logits/sampler and complete later native KV were not captured. A fresh prefill of a common prefix would not replace these missing live states.

\paragraph{Validation and retained evidence.}
The bridge's independent parser reconstructs 171 selected raw records and checks 1,596 per-run bindings; supplementary settings checks cover 183 engineering/formal records and 122 handoffs. It also reparses 114 historical native logs. All 80 local asset bindings pass post-execution readback. Its durable archive contains 2,276 logical files; full restoration and repeated CPU verification are inherited receipts, not newly repeated hardware execution in this revision. New document checks recompute the compact exports and retain all raw-scope limitations. Omitted model/runtime binaries and absolute execution time are not independently authenticated by content hashes.

\section{Offline Intervention-Selection Feasibility}
\label{app:selection}
\paragraph{Objective and action coverage.}
The primary decision is one configuration-level action restoring every baseline deviation while retaining all matching controls. Reference fallback is explicit and may avoid diagnosis entirely; retaining speculation would require a separate operational justification. A per-task mixture selected using known answers is not this policy. RTX19 supplies native N1 and ordinary U2; B57024 supplies native N1 and a qualified T1 proxy. MBPP24 has no measured U2. Missing entries remain unknown, and separate platform interventions do not form a complete factorial.

The analysis directly recomputes 285 RTX19 records (171 fresh bridge plus 57 old reference and 57 old H1), 360 B570 records and 360 already-exposed MBPP records. These 1,005 analysis records are not 1,005 new runs and do not equal the differently composed original intervention denominator. The 144 untruncated B570 exports retain 60/63 post-contract IDs. Repetition stability and all-three task agreement are checked without treating repetitions as independent decisions.

\paragraph{Two candidates and their limits.}
The stage/state candidate would compare baseline and reference live state at aligned boundaries and choose an action at the earliest differing stage. Required state and repair-direction evidence are absent, so no classifier or platform-to-remedy lookup is implemented. An observed difference need not identify a unique cause or a useful intervention.

The implemented candidate is generic charged proxy screening followed by actual verification. A query first charges acquisition, then reveals an available whole-workload outcome. The policy receives ordered action handles and replies, not task/platform names, reference tokens or unqueried native labels. It requests fixed reference/baseline first, screens a native action through its retained proxy if available, and verifies every accepted action. A failed proxy can cause false rejection on an untested configuration; it is not reported as a measured native failure. The policy stops on verified success, abstains after failure, or reports unresolved unknown coverage.

RTX screening requests 57 H1 and 57 U2 records after 114 common reference/baseline records and abstains. Enumeration requests 57 N1 and 57 U2 in either order and also abstains. B570 screening requests 72 T1 and 72 N1 after 144 common records; direct enumeration needs only 72 N1. Both select the same action. MBPP supplies no repairs and two persistent deviations; U2 is unknown in all 72 cells. Its results are not pooled with the other workloads. There are two specified policies and no outcome-driven policy revisions; one implementation correction removes unpaid baseline labels from the reference query without changing decisions.

\paragraph{Work and break-even accounting.}
For each native bridge condition, 57 runs evaluate 15,405 target caller tokens and 18,813 draft caller tokens, including prefill. Target continuation has 1,710 calls evaluating 5,118 tokens. N0 observes 1,704 two-token and 1,710 one-token graph occurrences; N1 observes 5,118 one-token occurrences. These are not unique kernel launches or latency measurements. Each condition has 57 handoffs; serialized payload volume is not measured device traffic. Four source-bound historical H1 spots establish activity in both proposal and target workers. Comparable proxy/native/reference timing and detailed U2/B570 graph costs are unavailable.

With proxy acceptance probability $p$, one-action screening costs $A+P+pN+V_P$ beyond common observations, compared with $N+V_E$ for direct verification. It can win only if $A+P+(V_P-V_E)<(1-p)N$ at an equivalent success/risk contract. The observed sole accepted B570 action offers no avoided N. The rejected RTX candidate could offer a saving only with a sufficiently cheap proxy, not established here. Setup may amortize over valid independent uses, but identical-workload caching also helps enumeration. Parser time and historical cost ratios cannot price these interventions.

\paragraph{Validation and decision.}
Ten CPU tests cover duplicate/missing keys, endpoint/cap handling, unknown denominators, paid observations, reference-label separation, charged probes, identifier invariance and hidden unqueried-label invariance. They test implementation integrity, not held-out method accuracy. No new model/GPU run occurs. The analysis recommends stopping this new-method route on the present evidence; it neither proves future selectors impossible nor demonstrates a new repair, causal diagnosis or saving.

\begin{table*}[t]
\centering
\caption{Supported decisions across distinct studies. Interventions and bridge tests have no qualified performance result. Historical cost studies retain their own reference, cohort and timing boundary.}
\label{tab:decisions}
\begin{tabular}{@{}>{\raggedright\arraybackslash}p{1.1in}>{\raggedright\arraybackslash}p{2.5in}>{\raggedright\arraybackslash}p{2.80in}@{}}
\toprule
Study & Observed result & Decision and limit \\
\midrule
B570 intervention & Four restorations; 20 matching controls retained & Restoration on the exposed cohort; no unique cause or measured gain \\
RTX intervention & One restoration, one new deviation; reference total unchanged & Reject the recovery/no-new-regression candidate \\
RTX simple baseline & Ordinary $U_2=N_1$ on 19/19 tasks & Complex predictor not shown necessary for this output objective \\
Offline selection & No RTX repair; B570 proxy adds work to sole accepted action & No demonstrated selector advantage; future costs remain unmeasured \\
RTX historical cost & Qualified serialization slower on all 19 prompts & Exclude measured serialization as a latency optimization \\
358H integration & Overhead-only prompt-GM ceilings below one & Target work/progress must change; no architectural impossibility \\
M4 runtime & Guard slower than Always-Spec; small fixed-mode opportunity & Prefer the measured fixed mode for this aggregate workload \\
\bottomrule
\end{tabular}
\end{table*}

